\chapter{Raízes Unitárias II}
\label{cap:raiz_unitaria2}
% ── índice ──────────────────────────────────────────────────────
\index{Phillips-Perron, teste de|textbf}
\index{phillips_perron@\texttt{phillips\_perron()}|textbf}
\index{Zivot-Andrews, teste de|textbf}
\index{zivot@\texttt{zivot()}|textbf}
\index{quebra estrutural|textbf}
\index{Newey-West, variância de|textbf}

% ─────────────────────────────────────────────────────────────────────────────
\section{Além do ADF: PP e Zivot-Andrews}

O capítulo anterior (cap.~\ref{cap:arima}) cobriu o teste ADF, que corrige
para autocorrelação residual por meio de defasagens. Dois problemas permanecem:

\begin{enumerate}
  \item O ADF é sensível à heterocedasticidade condicional dos erros (GARCH etc.)
  \item O ADF tem baixo poder na presença de \emph{quebra estrutural}: uma série
        estacionária com quebra de nível se parece com um passeio aleatório.
\end{enumerate}

\subsection*{Teste Phillips-Perron (PP)}

Phillips e Perron (1988) adotam uma abordagem não-paramétrica: corrigem a
estatística ADF para heterocedasticidade e autocorrelação de \emph{qualquer
forma} via estimador de Newey-West:

\[
  Z_\alpha = T(\hat{\rho} - 1) - \frac{T^2 \hat{\sigma}^2}{2\hat{s}_T^2}(\hat{\lambda}^2 - \hat{\gamma}_0)
\]

onde $\hat{s}_T^2$ é a variância de longo prazo estimada com kernel de Bartlett.
A distribuição assintótica é idêntica à do ADF (Dickey-Fuller).

\subsection*{Teste Zivot-Andrews (ZA)}

Zivot e Andrews (1992) relaxam a hipótese de ausência de quebra estrutural.
O teste estima a data de quebra \emph{endogenamente} como o valor que minimiza
a estatística $t$:

\[
  \hat{t}_B = \arg\min_\lambda \; t_\alpha(\lambda), \quad \lambda = t_B/T \in (0.15, 0.85)
\]

A equação testada inclui um nível de degrau ou mudança de tendência:
\[
  \Delta y_t = c + \beta t + \gamma D U_t(\lambda) + \alpha y_{t-1} + \sum_{j=1}^k d_j \Delta y_{t-j} + \varepsilon_t
\]

$H_0$: raiz unitária sem quebra estrutural.

% ─────────────────────────────────────────────────────────────────────────────
\section{Verificação por DGP}

\begin{lstlisting}[caption={PP e Zivot-Andrews}]
seed(42)
// Passeio aleatório: rw_t = rw_{t-1} + e_t
// Série estacionária com tendência: stat_t = 0.5*t + e
// Série com quebra em t=150: AR(1) + degrau de 3.0

phillips_perron(df, rw)
phillips_perron(df, stat)
zivot(df, break_t)
\end{lstlisting}

\begin{lstlisting}[language={}, basicstyle=\ttfamily\small, frame=none,
  numbers=none, backgroundcolor=\color{codbg}]
PP — rw:      Zt = 77.365   VC 5%=-2.860   NÃO rejeita H₀  (raiz unitária) ✓
PP — stat:    Zt = 2058     VC 5%=-2.860   NÃO rejeita H₀  (trend-stationary)
ZA — break_t: t  = -11.597  VC 5%=-4.800   REJEITA H₀, quebra em obs=147 ✓
\end{lstlisting}

O PP para \texttt{stat} (tendência linear + ruído) não rejeita $H_0$ porque o
teste padrão (sem tendência determinística na equação auxiliar) não distingue
tendência estocástica de determinística. O ZA identifica corretamente a quebra
estrutural em $t=147$ (próximo ao verdadeiro $t=150$) e rejeita $H_0$.

% ─────────────────────────────────────────────────────────────────────────────
\section{Sintaxe}

\begin{lstlisting}
phillips_perron(df, var)         // PP sem tendência
phillips_perron(df, var, trend=true)  // PP com tendência determinística
zivot(df, var)                   // Zivot-Andrews (quebra endógena de nível)
zivot(df, var, type="trend")     // quebra na tendência
\end{lstlisting}

\begin{center}
\begin{tabular}{llll}
\toprule
\textbf{Teste} & \textbf{Correção} & \textbf{Quebra?} & \textbf{VC 5\%} \\
\midrule
ADF      & Defasagens paramétricas & Não & $-2{,}86$ \\
PP       & Newey-West não-param.   & Não & $-2{,}86$ \\
ZA       & PP + quebra endógena    & Sim & $-4{,}80$ \\
\bottomrule
\end{tabular}
\end{center}

\begin{nota}
  Séries trend-stationary e séries com quebra estrutural são dois casos onde
  o ADF/PP pode levar a conclusões erradas. O ZA resolve o segundo; para o
  primeiro, incluir uma tendência determinística na equação auxiliar
  (\texttt{trend=true}) é suficiente.
\end{nota}
